% ============================================================================
% AI/ML LECTURE NOTES - PREAMBLE
% Custom LaTeX setup for comprehensive course notes
% ============================================================================

% ---- Document Class ----
% Using report for chapter structure, adjust font size as needed
\documentclass[11pt,a4paper,twoside]{report}

% ============================================================================
% ESSENTIAL PACKAGES
% ============================================================================

% ---- Mathematics ----
\usepackage{amsmath,amssymb,amsthm}
\usepackage{mathtools}
\usepackage{bm}                    % Bold math symbols

% ---- Fonts and Typography ----
\usepackage[utf8]{inputenc}
\usepackage[T1]{fontenc}
\usepackage{lmodern}
\usepackage{microtype}             % Better typography

% ---- Layout ----
\usepackage[margin=2.5cm]{geometry}
\usepackage{fancyhdr}
\usepackage{titlesec}
\usepackage{parskip}               % Space between paragraphs

% ---- Graphics and Diagrams ----
\usepackage{graphicx}
\usepackage{tikz}
\usepackage{pgfplots}
\pgfplotsset{compat=1.18}
\usetikzlibrary{arrows.meta, positioning, shapes.geometric, calc, decorations.pathreplacing}

% ---- Tables ----
\usepackage{booktabs}
\usepackage{array}
\usepackage{multirow}
\usepackage{longtable}

% ---- Colors and Boxes ----
\usepackage[dvipsnames,svgnames,x11names]{xcolor}
\usepackage{tcolorbox}
\tcbuselibrary{breakable,skins,theorems}

% ---- Code Listings ----
\usepackage{listings}

% ---- Hyperlinks ----
\usepackage{hyperref}
\hypersetup{
    colorlinks=true,
    linkcolor=NavyBlue,
    citecolor=ForestGreen,
    urlcolor=RoyalBlue,
    pdfauthor={AI/ML Course},
    pdftitle={Machine Learning and Deep Learning Notes}
}

% ---- Bibliography (optional) ----
\usepackage{natbib}

% ============================================================================
% CUSTOM COLORS
% ============================================================================

\definecolor{additionalblue}{RGB}{230, 242, 255}
\definecolor{additionalborder}{RGB}{66, 133, 244}
\definecolor{definitiongreen}{RGB}{232, 245, 233}
\definecolor{definitionborder}{RGB}{76, 175, 80}
\definecolor{theoremyellow}{RGB}{255, 249, 222}
\definecolor{theoremborder}{RGB}{255, 193, 7}
\definecolor{examplegray}{RGB}{245, 245, 245}
\definecolor{exampleborder}{RGB}{158, 158, 158}
\definecolor{codebg}{RGB}{248, 248, 248}
\definecolor{keywordcolor}{RGB}{0, 100, 180}
\definecolor{stringcolor}{RGB}{163, 21, 21}
\definecolor{commentcolor}{RGB}{0, 128, 0}

% ============================================================================
% CUSTOM THEOREM ENVIRONMENTS
% ============================================================================

\theoremstyle{definition}
\newtheorem{definition}{Definition}[chapter]
\newtheorem{example}{Example}[chapter]

\theoremstyle{plain}
\newtheorem{theorem}{Theorem}[chapter]
\newtheorem{lemma}[theorem]{Lemma}
\newtheorem{corollary}[theorem]{Corollary}
\newtheorem{proposition}[theorem]{Proposition}

\theoremstyle{remark}
\newtheorem*{remark}{Remark}
\newtheorem*{note}{Note}

% ============================================================================
% CUSTOM TCOLORBOX ENVIRONMENTS
% ============================================================================

% ---- Additional Context (Blue Box) - For extra explanations ----
\newtcolorbox{additionalcontext}[1][Additional Context]{
    enhanced,
    breakable,
    colback=additionalblue,
    colframe=additionalborder,
    coltitle=white,
    fonttitle=\bfseries,
    title={\faLightbulbO\ #1},
    before upper={\parindent0pt},
    boxrule=1pt,
    arc=3pt,
    left=10pt,
    right=10pt,
    top=8pt,
    bottom=8pt,
    before skip=12pt,
    after skip=12pt,
    attach boxed title to top left={yshift=-2mm, xshift=5mm},
    boxed title style={
        colback=additionalborder,
        arc=2pt,
        boxrule=0pt
    }
}

% Simplified version without icon (in case fontawesome not available)
\newtcolorbox{bluebox}[1][Additional Context]{
    enhanced,
    breakable,
    colback=additionalblue,
    colframe=additionalborder,
    coltitle=white,
    fonttitle=\bfseries,
    title={#1},
    before upper={\parindent0pt},
    boxrule=1.5pt,
    arc=4pt,
    left=10pt,
    right=10pt,
    top=10pt,
    bottom=10pt,
    before skip=15pt,
    after skip=15pt,
}

% ---- Definition Box ----
\newtcolorbox{definitionbox}[1][Definition]{
    enhanced,
    breakable,
    colback=definitiongreen,
    colframe=definitionborder,
    coltitle=black,
    fonttitle=\bfseries,
    title={#1},
    boxrule=1pt,
    arc=3pt,
    before skip=12pt,
    after skip=12pt,
}

% ---- Theorem Box ----
\newtcolorbox{theorembox}[1][Theorem]{
    enhanced,
    breakable,
    colback=theoremyellow,
    colframe=theoremborder,
    coltitle=black,
    fonttitle=\bfseries,
    title={#1},
    boxrule=1pt,
    arc=3pt,
    before skip=12pt,
    after skip=12pt,
}

% ---- Example Box ----
\newtcolorbox{examplebox}[1][Example]{
    enhanced,
    breakable,
    colback=examplegray,
    colframe=exampleborder,
    coltitle=black,
    fonttitle=\bfseries,
    title={#1},
    boxrule=0.5pt,
    arc=2pt,
    before skip=10pt,
    after skip=10pt,
}

% ---- Important Note Box ----
\newtcolorbox{importantbox}[1][Important]{
    enhanced,
    breakable,
    colback=red!5,
    colframe=red!70!black,
    coltitle=white,
    fonttitle=\bfseries,
    title={#1},
    boxrule=1.5pt,
    arc=3pt,
    before skip=12pt,
    after skip=12pt,
}

% ---- Key Concept Box (Blue with emphasis) ----
\newtcolorbox{keyconcept}{
    enhanced,
    breakable,
    colback=blue!8,
    colframe=blue!70!black,
    coltitle=white,
    fonttitle=\bfseries,
    title={Key Concept},
    boxrule=1.5pt,
    arc=4pt,
    before skip=15pt,
    after skip=15pt,
}

% ---- Important Highlight (Red/Orange) ----
\newtcolorbox{important}{
    enhanced,
    breakable,
    colback=orange!10,
    colframe=orange!80!black,
    boxrule=1.5pt,
    arc=3pt,
    left=10pt,
    right=10pt,
    top=8pt,
    bottom=8pt,
    before skip=12pt,
    after skip=12pt,
}

% ---- Info Box (Gray informational) ----
\newtcolorbox{infobox}[1][Information]{
    enhanced,
    breakable,
    colback=gray!8,
    colframe=gray!60,
    coltitle=black,
    fonttitle=\bfseries,
    title={#1},
    boxrule=1pt,
    arc=3pt,
    before skip=12pt,
    after skip=12pt,
}

% ---- Derivation Box (Purple for mathematical derivations) ----
\newtcolorbox{derivation}[1][Derivation]{
    enhanced,
    breakable,
    colback=purple!5,
    colframe=purple!60!black,
    coltitle=white,
    fonttitle=\bfseries,
    title={#1},
    boxrule=1pt,
    arc=3pt,
    before skip=12pt,
    after skip=12pt,
}

% ---- Algorithm Box (Teal for algorithms) ----
\newtcolorbox{algorithm}[1][Algorithm]{
    enhanced,
    breakable,
    colback=teal!8,
    colframe=teal!70!black,
    coltitle=white,
    fonttitle=\bfseries,
    title={#1},
    boxrule=1pt,
    arc=3pt,
    before skip=12pt,
    after skip=12pt,
}

% ---- Summary Box (Green for chapter summaries) ----
\newtcolorbox{summary}{
    enhanced,
    breakable,
    colback=green!8,
    colframe=green!60!black,
    coltitle=white,
    fonttitle=\bfseries,
    title={Summary},
    boxrule=1.5pt,
    arc=4pt,
    before skip=15pt,
    after skip=15pt,
}

% ---- Simple formulabox environment (for formulas) ----
\newenvironment{formulabox}{%
    \begin{center}
    \begin{tcolorbox}[
        enhanced,
        colback=gray!5,
        colframe=gray!50,
        boxrule=0.5pt,
        arc=2pt,
        width=0.9\textwidth,
        left=10pt,
        right=10pt,
        top=8pt,
        bottom=8pt,
    ]
}{%
    \end{tcolorbox}
    \end{center}
}

% ============================================================================
% CODE LISTING SETTINGS
% ============================================================================

\lstdefinestyle{pythonstyle}{
    language=Python,
    backgroundcolor=\color{codebg},
    basicstyle=\ttfamily\small,
    keywordstyle=\color{keywordcolor}\bfseries,
    stringstyle=\color{stringcolor},
    commentstyle=\color{commentcolor}\itshape,
    numberstyle=\tiny\color{gray},
    numbers=left,
    numbersep=8pt,
    stepnumber=1,
    showstringspaces=false,
    breaklines=true,
    breakatwhitespace=true,
    frame=single,
    framerule=0.5pt,
    rulecolor=\color{gray!50},
    tabsize=4,
    captionpos=b,
    morekeywords={self, True, False, None, as, with, yield},
    deletekeywords={abs, min, max, sum, len, range, print, input, open, list, dict, set, tuple, str, int, float, bool},
    xleftmargin=15pt,
    xrightmargin=5pt,
    aboveskip=12pt,
    belowskip=12pt,
}

\lstset{style=pythonstyle}

% Python code environment
\lstnewenvironment{pythoncode}[1][]{
    \lstset{style=pythonstyle, #1}
}{}

% ============================================================================
% MATHEMATICAL NOTATION COMMANDS
% ============================================================================

% ---- Vectors and Matrices ----
\newcommand{\vect}[1]{\mathbf{#1}}           % Vectors: \vect{x}
\newcommand{\mat}[1]{\mathbf{#1}}            % Matrices: \mat{W}
\newcommand{\transpose}{^\top}                % Transpose: \transpose

% ---- Derivatives ----
\newcommand{\deriv}[2]{\frac{d #1}{d #2}}                    % d/dx
\newcommand{\pderiv}[2]{\frac{\partial #1}{\partial #2}}     % Partial derivative
\newcommand{\grad}{\nabla}                                    % Gradient symbol

% ---- Probability and Statistics ----
\newcommand{\expect}[1]{\mathbb{E}\left[#1\right]}           % Expectation
\newcommand{\prob}[1]{\mathbb{P}\left(#1\right)}             % Probability
\newcommand{\var}[1]{\text{Var}\left(#1\right)}              % Variance
\newcommand{\cov}[2]{\text{Cov}\left(#1, #2\right)}          % Covariance

% ---- Sets and Spaces ----
\newcommand{\R}{\mathbb{R}}                  % Real numbers
\newcommand{\N}{\mathbb{N}}                  % Natural numbers
\newcommand{\Z}{\mathbb{Z}}                  % Integers

% ---- Neural Network Notation ----
\newcommand{\sigmoid}[1]{\sigma\left(#1\right)}              % Sigmoid
\newcommand{\relu}[1]{\text{ReLU}\left(#1\right)}            % ReLU
\newcommand{\softmax}[1]{\text{softmax}\left(#1\right)}      % Softmax
\newcommand{\loss}{\mathcal{L}}                               % Loss function
\newcommand{\pred}[1]{\hat{#1}}                               % Prediction hat

% ---- Common Operations ----
\newcommand{\argmin}{\mathop{\mathrm{arg\,min}}}
\newcommand{\argmax}{\mathop{\mathrm{arg\,max}}}

\pagestyle{fancy}
\fancyhf{}

% Custom command to set the current lecture tracking
\newcommand{\setlecture}[1]{%
    \def\currentlecture{Lecture #1}%
}
% Default value if not set
\def\currentlecture{}

\fancyhead[LE,RO]{\thepage}
\fancyhead[LO]{\nouppercase{\rightmark} \ifx\currentlecture\empty\else\quad|\quad \textbf{\currentlecture}\fi}
\fancyhead[RE]{\nouppercase{\leftmark} \ifx\currentlecture\empty\else\textbf{\currentlecture} \quad|\quad \fi}
\renewcommand{\headrulewidth}{0.4pt}

% ---- Chapter Title Formatting ----
\titleformat{\chapter}[display]
    {\normalfont\huge\bfseries}
    {\chaptertitlename\ \thechapter}
    {20pt}
    {\Huge}

\titlespacing*{\chapter}{0pt}{-30pt}{40pt}

% ============================================================================
% FIGURE PLACEHOLDER COMMAND
% ============================================================================

\newcommand{\figurePlaceholder}[2]{%
    \begin{figure}[htbp]
        \centering
        \begin{tikzpicture}
            \node[draw=gray, dashed, minimum width=0.8\textwidth, minimum height=4cm, 
                  align=center, text=gray] {
                \textbf{[FIGURE PLACEHOLDER]}\\[5pt]
                #2
            };
        \end{tikzpicture}
        \caption{#1}
    \end{figure}
}

% ============================================================================
% DIAGRAM PLACEHOLDER COMMAND
% ============================================================================

\newcommand{\diagramPlaceholder}[2]{%
    \begin{center}
        \begin{tikzpicture}
            \node[draw=blue!50, dashed, minimum width=0.7\textwidth, minimum height=3cm, 
                  align=center, fill=blue!5] {
                \textbf{[DIAGRAM: #1]}\\[5pt]
                \textit{File: #2}
            };
        \end{tikzpicture}
    \end{center}
}
